% This contents of this file will be inserted into the _Solutions version of the
% output tex document.  Here's an example:

% If assignment with subquestion (1.a) requires a written response, you will
% find the following flag within this document: <SCPD_SUBMISSION_TAG>_1a
% In this example, you would insert the LaTeX for your solution to (1.a) between
% the <SCPD_SUBMISSION_TAG>_1a flags.  If you also constrain your answer between the
% START_CODE_HERE and END_CODE_HERE flags, your LaTeX will be styled as a
% solution within the final document.

% Please do not use the '<SCPD_SUBMISSION_TAG>' character anywhere within your code.  As expected,
% that will confuse the regular expressions we use to identify your solution.

\def\assignmentnum{3 }
% \def\assignmentname{Sentiment Analysis}
\def\assignmenttitle{XCS234 Assignment \assignmentnum}
\input{macros}
\begin{document}
\pagestyle{myheadings} \markboth{}{\assignmenttitle}

% <SCPD_SUBMISSION_TAG>_entire_submission

This handout includes space for every question that requires a written response.
Please feel free to use it to handwrite your solutions (legibly, please).  If
you choose to typeset your solutions, the |README.md| for this assignment includes
instructions to regenerate this handout with your typeset \LaTeX{} solutions.
\ruleskip


\LARGE
2.a
\normalsize

% <SCPD_SUBMISSION_TAG>_2a
\begin{answer}
  % ### START CODE HERE ###
  The probability of a specific trajectory is the product of the probabilities of the agent choosing an action based on its policy and the environment transitioning to the next state based on that action.

  $$
  \rho^{\pi}(\tau) = p(s_0) \prod_{t=0}^{\infty} \pi(a_t|s_t) \mathcal{P}(s_{t+1}|s_t, a_t)
  $$

  Because we assume that the MDP has a single, fixed starting state $s_0$, i.e. $p(s_0) = 1$. Therefore,

  $$
  \rho^{\pi}(\tau) = \prod_{t=0}^{\infty} \pi(a_t|s_t) \mathcal{P}(s_{t+1}|s_t, a_t)
  $$
  % ### END CODE HERE ###
\end{answer}
% <SCPD_SUBMISSION_TAG>_2a
\clearpage

\LARGE
2.b
\normalsize

% <SCPD_SUBMISSION_TAG>_2b
\begin{answer}
  % ### START CODE HERE ###
  Firstly, apply the linearity of expectation and then expand the expectation.

  \begin{align*}
  \mathbb{E}_{\tau\sim\rho^{\pi}}\left[\sum_{t=0}^{\infty}\gamma^{t}f(s_{t},a_{t})\right] &= \sum_{t=0}^{\infty}\gamma^{t}\mathbb{E}_{\tau\sim\rho^{\pi}}[f(s_{t},a_{t})] \\
    &= \sum_{t=0}^{\infty}\gamma^{t}\sum_{s \in \mathcal{S}} \sum_{a \in \mathcal{A}} p(s_t=s, a_t=a) f(s,a) \\
    &= \sum_{t=0}^{\infty}\gamma^{t}\sum_{s \in \mathcal{S}} \sum_{a \in \mathcal{A}} p(s_t=s) p(a_t=a | s_t=s) f(s,a) \\
    &= \sum_{t=0}^{\infty}\gamma^{t}\sum_{s \in \mathcal{S}} \sum_{a \in \mathcal{A}} p(s_t=s) \pi(a|s) f(s,a) \\
    &= \sum_{s \in \mathcal{S}} \sum_{a \in \mathcal{A}} \pi(a|s) f(s,a) \left(\sum_{t=0}^{\infty}\gamma^{t} p(s_t=s)\right)
  \end{align*}

  We have the definition of the discounted, stationary state distribution

  $$
  d^{\pi}(s) = (1-\gamma)\sum_{t=0}^{\infty}\gamma^{t}p(s_{t}=s)
  $$


  Rearrange the equation we can get solve $\sum_{t=0}^{\infty}\gamma^{t} p(s_t=s)$

  $$
  \sum_{t=0}^{\infty}\gamma^{t}p(s_t=s) = \frac{d^{\pi}(s)}{1-\gamma}
  $$

  Substitute back and we get

  \begin{align*}
  \mathbb{E}_{\tau\sim\rho^{\pi}}\left[\sum_{t=0}^{\infty}\gamma^{t}f(s_{t},a_{t})\right] &= \sum_{s \in \mathcal{S}} \sum_{a \in \mathcal{A}} \pi(a|s) f(s,a) \left(\frac{d^{\pi}(s)}{1-\gamma}\right) \\
    &= \frac{1}{1-\gamma} \sum_{s \in \mathcal{S}} d^{\pi}(s) \left(\sum_{a \in \mathcal{A}} \pi(a|s) f(s,a)\right) \\
    &= \frac{1}{1-\gamma}\mathbb{E}_{s\sim d^{\pi}}[\mathbb{E}_{a\sim\pi(s)}[f(s,a)]]
  \end{align*}

  % ### END CODE HERE ###
\end{answer}
% <SCPD_SUBMISSION_TAG>_2b
\clearpage

\LARGE
2.c
\normalsize

% <SCPD_SUBMISSION_TAG>_2c
\begin{answer}
  \begin{align*}
      V^\pi(s_0) - V^{\pi'}(s_0) &= \mathbb{E}_{\tau \sim \rho^\pi}\left[\sum\limits_{t=0}^\infty \gamma^t \mathcal{R}(s_t,a_t)\right] - V^{\pi'}(s_0) \\
      &= \mathbb{E}_{\tau \sim \rho^\pi}\left[\sum\limits_{t=0}^\infty \gamma^t\left( \mathcal{R}(s_t,a_t) + V^{\pi'}(s_t) - V^{\pi'}(s_t)\right)\right] - V^{\pi'}(s_0) \\
      &= \mathbb{E}_{\tau \sim \rho^\pi}\left[\sum\limits_{t=0}^\infty \gamma^t\left( \mathcal{R}(s_t,a_t) + \gamma V^{\pi'}(s_{t+1}) - V^{\pi'}(s_t)\right)\right] \\
  % ### START CODE HERE ###
  &= \mathbb{E}_{\tau \sim \rho^\pi}\left[\sum\limits_{t=0}^\infty \gamma^t\left( Q^{\pi'}(s_t,a_t) - V^{\pi'}(s_t)\right)\right] && \text{Definition of $Q^{\pi'}(s_t,a_t)$} \\
  &= \mathbb{E}_{\tau \sim \rho^\pi}\left[\sum\limits_{t=0}^\infty \gamma^t A^{\pi'}(s_t,a_t)\right] && \text{Definition of $A^{\pi'}(s_t,a_t)$} \\
  &= \frac{1}{(1-\gamma)}\mathbb{E}_{s\sim d^{\pi}}[\mathbb{E}_{a\sim\pi(s)}[A^{\pi^{\prime}}(s,a)]] && \text{From 2.b}
  % ### END CODE HERE ###
  \end{align*}
  \end{answer}
% <SCPD_SUBMISSION_TAG>_2c
\clearpage

\LARGE
3.a
\normalsize

% <SCPD_SUBMISSION_TAG>_3a
\begin{answer}
  % ### START CODE HERE ###
  1. To follow *Respect for persons*, the researcher must properly
  inform the students about the study's purpose, the method they
  will implement including the use of AI chatbot and random assignment
  and any other possible risk with the student's consent.

  2. To minimize possible harms (as stated in *Beneficence*), the researcher
  should ensure that the students can withdraw from the study at any time
  without penalty if they feel the chatbot is ineffective.
  % ### END CODE HERE ###
\end{answer}
% <SCPD_SUBMISSION_TAG>_3a
\clearpage

% <SCPD_SUBMISSION_TAG>_entire_submission

\end{document}
